\documentclass[11pt]{article}
\usepackage[a4paper,margin=1in]{geometry}
\usepackage{amsmath,amssymb,amsthm,mathtools}
\usepackage{booktabs}
\usepackage{enumitem}
\usepackage{microtype}
\usepackage{hyperref}
\usepackage{xcolor}
\hypersetup{colorlinks=true,linkcolor=blue!50!black,citecolor=blue!50!black,urlcolor=blue!50!black}

\newtheorem{theorem}{Theorem}[section]
\newtheorem{proposition}[theorem]{Proposition}
\newtheorem{lemma}[theorem]{Lemma}
\newtheorem{corollary}[theorem]{Corollary}
\theoremstyle{definition}
\newtheorem{definition}[theorem]{Definition}
\theoremstyle{remark}
\newtheorem{remark}[theorem]{Remark}
\newcommand{\J}{J}
\newcommand{\Bzero}{B_0}
\newcommand{\Bone}{B_1}
\newcommand{\Leb}{\operatorname{Leb}}
\newcommand{\spec}{\operatorname{spec}}

\title{Boundary-Aligned Cyclic-Ulam Structure and Phase-Local Leakage\\
\large An exact finite-dimensional edge at a postcritically finite quadratic map}
\author{RH research program}
\date{August 2026}

\begin{document}
\maketitle

\begin{abstract}
We isolate an exact finite-dimensional mechanism in a postcritically finite
quadratic map.  At the algebraic parameter $u_c$ satisfying
$u_c^3-2u_c^2+2u_c-2=0$, the invariant interval splits into two bands that
are exchanged by one iterate.  For every finite exact cell-overlap Ulam
partition whose boundary contains the band endpoint $r=u_c-1$, the matrix is
anti-diagonal in band coordinates and inherits a sign vector with eigenvalue
$-1$.  If one cell crosses the endpoint, the projected sign has the exact
local defect $4\theta(1-\theta)$, or $4h\theta(1-\theta)$ after weighting by
the cell width.  A frozen 33-phase scan at four resolutions reproduces the
predicted aligned/crossing dichotomy and records the finite near-$-1$ drift.
The scan is a diagnostic: no universal power law, continuum isolated-spectrum
theorem, canonical arithmetic operator, or Riemann-Hypothesis implication is
claimed.  An explicit overlap ledger separates this edge from RH-3, RH-10,
and RH-55.
\end{abstract}

\section{Scope and claim boundary}

The object of this paper is a finite exact Ulam matrix, not an unnamed
continuum transfer operator.  We use the phrase ``cyclic mode'' for the
two-band exchange and ``leakage'' for a projected finite-cell diagnostic.
The following statements are outside the paper's claim ceiling: a
strong-Banach-space isolated eigenvalue or a common resolvent contour as the
mesh is refined; a universal $\sqrt{\sigma}$ law (or any exponent inferred
from a finite log--log fit); and a canonical global H\'enon operator, zeta
factor, von Mangoldt trace, Hilbert--Polya construction, Riemann-zero
identification, or proof of RH.
All numerical rows are frozen source diagnostics and are not asymptotic
evidence.  The exact route verdict is Route A $=\texttt{GO}$ and Route B
$=\texttt{STOP\_SCOPED}$.

\section{The postcritically finite two-band map}

Let
\begin{equation}
 p(u)=u^3-2u^2+2u-2,
 \qquad u_c=1.5436890126920764\ldots,
 \qquad r=u_c-1.
\end{equation}
The map and its invariant interval are
\begin{equation}
 f(x)=1-u_cx^2,
 \qquad \J=[-r,1].
\end{equation}
The symbolic source certificate records
\begin{equation}
 f(0)=1,\qquad f(1)=-r,\qquad f(-r)=r,\qquad f(r)=r.
 \label{eq:critical}
\end{equation}
Define the two closed bands (with their common endpoint understood in the
usual measure-zero convention)
\begin{equation}
 \Bzero=[-r,r],\qquad \Bone=[r,1].
\end{equation}
Monotonicity on either side of zero and \eqref{eq:critical} give
\begin{equation}
 f(\Bzero)=\Bone,\qquad f(\Bone)=\Bzero.
 \label{eq:exchange}
\end{equation}

The source paper also records the normalized invariant density $h$ and the
$L^1$ sign mode
\begin{equation}
 g=h(\mathbf 1_{\Bzero}-\mathbf 1_{\Bone}),
 \qquad \mathcal P g=-g,
 \label{eq:l1mode}
\end{equation}
in the invariant-density normalization.  Equation \eqref{eq:l1mode} is an
$L^1$ statement.  It does not assert that $-1$ is isolated on a selected
strong space.

\section{Exact aligned Ulam inheritance}

\begin{definition}[Exact cell-overlap matrix]
Let $\mathcal D_h=\{D_i\}_{i=1}^N$ be a finite interval partition of $\J$.
The row-stochastic exact cell-overlap matrix is
\begin{equation}
 (P_h)_{ij}=\frac{\Leb(D_i\cap f^{-1}(D_j))}{\Leb(D_i)}.
 \label{eq:ulam}
\end{equation}
We call the partition aligned when $r$ is a cell boundary.  Source cells are
split at the turning point and inverse-branch breakpoints before the overlap
integrals are evaluated.
\end{definition}

\begin{theorem}[Aligned block and sign mode]\label{thm:block}
For every aligned finite partition, after ordering cells by band,
\begin{equation}
 P_h=\begin{pmatrix}0&A\\B&0\end{pmatrix}.
 \label{eq:block}
\end{equation}
If $s_i=+1$ on $\Bzero$ cells and $s_i=-1$ on $\Bone$ cells, then
\begin{equation}
 P_hs=-s.
 \label{eq:sign}
\end{equation}
In particular, $-1\in\spec(P_h)$ in exact arithmetic.
\end{theorem}

\begin{proof}
Every source cell belongs to one band.  By \eqref{eq:exchange}, its image
belongs to the opposite band, so an overlap between cells in the same band is
zero.  This is exactly the anti-diagonal form \eqref{eq:block}.  Every
nonzero row therefore sends its mass to cells carrying the opposite sign;
multiplication by $P_h$ changes $s$ to $-s$, which proves \eqref{eq:sign}.
No uniform mesh assumption is used.
\end{proof}

\begin{remark}
The theorem is finite-dimensional and normalization-specific.  It does not
identify the vector $s$ with a right eigenfunction of an arbitrary Lebesgue
reference Perron--Frobenius operator, and it does not control the rest of the
finite spectrum or its limit as $h\to0$.
\end{remark}

\section{Crossing-cell projection defect}

Suppose a single cell $D_*$ crosses $r$.  Define the band fraction and the
projected sign by
\begin{equation}
 q_i=\frac{\Leb(D_i\cap\Bzero)}{\Leb(D_i)},
 \qquad s_i^{(h)}=2q_i-1.
 \label{eq:q}
\end{equation}
For an aligned cell, $q_i\in\{0,1\}$.  For the crossing cell,
$q_* =\theta\in(0,1)$.  The elementary identity below is the exact local
phase law.

\begin{proposition}[Local phase defect]\label{prop:defect}
For $\theta\in[0,1]$,
\begin{equation}
 1-(2\theta-1)^2=4\theta(1-\theta).
 \label{eq:identity}
\end{equation}
If the crossing cell has width $h$, its width-weighted defect is
$4h\theta(1-\theta)$.
\end{proposition}

\begin{proof}
Expand $(2\theta-1)^2=4\theta^2-4\theta+1$ and subtract from one.
Multiplying by $h$ gives the second assertion.
\end{proof}

The global diagnostic used by the source is the stationary same-band mass
\begin{equation}
 L_h=\sum_i\pi_iq_i\sum_j(P_h)_{ij}q_j
 +\sum_i\pi_i(1-q_i)\sum_j(P_h)_{ij}(1-q_j),
 \label{eq:global}
\end{equation}
where $\pi$ is a stationary cell-mass vector.  If the partition is aligned,
the anti-diagonal theorem makes $L_h=0$ exactly.  When a cell crosses $r$,
$L_h$ depends on the stationary weights and all other cells; it is not equal
to the local expression in every discretization.  This distinction matters:
the local identity is a theorem, while the global phase scan is an observed
finite quantity.

\section{Frozen numerical protocol}

The external source uses analytic inverse-branch overlap assembly, row
normalization, stationary iteration, and a tracked eigenvalue nearest $-1$.
The retained deterministic scan has four resolutions
\[
 N\in\{256,512,1024,2048\}
\]
and 33 translated phases at each resolution, together with one snapped
aligned phase.  Thus there are 136 rows: 4 aligned and 132 crossing.  The
aligned rows have zero projected same-band mass and an exact sign residual up
to floating assembly error.  Crossing rows have positive projected mass and a
nonzero finite displacement of the tracked near-$-1$ eigenvalue.

\begin{table}[t]
\centering
\caption{Frozen phase-scan counts and ranges.  These values are diagnostics.}
\label{tab:scan}
\begin{tabular}{@{}lr@{}}
\toprule
quantity & value\\
\midrule
resolutions & $256,512,1024,2048$\\
phases per resolution & $33$\\
total rows & $136$\\
aligned / crossing rows & $4/132$\\
crossing $L_h$ range & $1.14\times10^{-4}$ to $7.71\times10^{-3}$\\
maximum crossing distance from $-1$ & $4.91\times10^{-3}$\\
maximum row-sum error & $6.4\times10^{-9}$\\
maximum stationary $L^1$ residual & $1.9\times10^{-13}$\\
\bottomrule
\end{tabular}
\end{table}

The source also contains direct, split, and cycle-constrained noisy matrices.
Their fitted slopes over a finite range are deliberately not promoted here.
The boundary convention and kernel change the finite intercepts, and no
common strong-space perturbation theorem is supplied.

\section{Overlap audit and route decision}

RH-3 uses the same band-merging map to prove a continuum parity eigenmode,
periodogram limits, and conditional two-step operator decompositions.  It
does not contain the arbitrary aligned finite-Ulam block theorem or the
crossing-cell identity.  RH-10 studies exact periodic counts, exponentially
close boundary cycles, noncommuting long-cycle/noise limits, and a
parity-renormalized determinant; it does not state the phase-local Ulam
defect.  RH-55 proves midpoint--Ulam strong--weak contour transfer for a
conditioned folded-Gaussian kernel; it does not prove deterministic PCF sign
inheritance for arbitrary aligned partitions.  The overlap is therefore
structural context, not duplication.

Route A is \texttt{GO} for Theorems~\ref{thm:block} and
\ref{prop:defect}, together with the source-locked finite protocol.  Route B
is \texttt{STOP\_SCOPED}: the first missing bridge is a common strong-space
projector/resolvent theorem connecting the finite matrices to a continuum
operator.  In particular, no finite scan proves a universal noise exponent,
an isolated continuum resonance, or any arithmetic/RH statement.

\section{Conclusion}

The exact content of RH-367 is a reusable finite-volume fact: a genuine
band-boundary alignment preserves a two-cycle sign mode in every exact
cell-overlap matrix, while a crossing cell has an explicit phase defect.  The
global leakage and near-$-1$ displacement are useful diagnostics for choosing
what must be proved next, but they are not substitutes for a strong-space
limit theorem.  The physical route coordinate remains
\texttt{actual\_same\_clock\_unnormalized\_head\_transport\_open}; all five
RH gates remain false/open.

\begingroup\raggedright
\begin{thebibliography}{9}
\bibitem{cyclic}
  Cyclic Ulam Map, source package and proof bundle, 2026.  Frozen commit
  \texttt{e7d21f646498d77e1c3213d1e4f35dc8466038ff}.
\bibitem{rh3}
  RH-3, \emph{Parity-Resolved Dynamics at a Quadratic Band-Merging
  Parameter}, repository artifact.
\bibitem{rh10}
  RH-10, \emph{Parity-Renormalized Long-Cycle Determinants}, repository
  artifact.
\bibitem{rh55}
  RH-55, \emph{Strong--Weak Adaptive Cutoff Transfer for Intrinsic Riesz
  Factors}, repository artifact.
\bibitem{ulam}
  G. Froyland, ``On Ulam approximation of isolated spectrum and
  eigenfunctions of hyperbolic maps,'' \emph{DCDS} 17 (2007), 671--689.
\end{thebibliography}
\endgroup

\end{document}
